\chapter{Requirements and System Design}
\label{ch:requirements-system-design}

The previous chapter established the technical foundations and related-work
context for the proposed monitoring system. This chapter translates that
background into explicit requirements and architectural decisions. Its purpose
is to explain what the system is expected to do, which constraints shape its
design, and why responsibilities are divided between kernel space and user
space.

The design follows one central principle: kernel programs should perform only
the bounded work that benefits directly from proximity to the observed
operation, while the more changeable interpretation and presentation logic
should remain in user space. An explicit event contract connects these two
domains. Policies then constrain the admitted event stream, and independently
loaded detectors evaluate either individual observations or short ordered
sequences. This progression connects the requirements to the architecture
without anticipating the implementation details presented in Chapter~4 or the
experimental evidence presented in Chapter~5.

\section{Design Goals and Constraints}
\label{sec:design-goals-constraints}

The system is intended to provide focused host-level visibility rather than an
exhaustive account of all kernel activity. Its event domain consequently
centres on process execution and security-relevant operations involving
credentials, files, memory, namespaces, capabilities, and kernel mechanisms.
The architecture must make this domain selectable, convert heterogeneous
kernel observations into one typed event model, and support deterministic
analysis without placing the entire detection model inside the eBPF programs.

Three concerns influence every stage of the design. First, compatibility must
be established on the concrete enterprise kernel rather than inferred from an
upstream version. Second, high-frequency observations require early scope
reduction and bounded state. Third, failures at the producer, transport,
decoder, or detector boundary must remain visible so that missing telemetry is
not mistaken for an absence of monitored activity.

\subsection{Target Environment and Compatibility Assumptions}
\label{subsec:target-environment-assumptions}

The reference environment is Rocky Linux 8.10 on x86-64. The exact kernel
build is \path{4.18.0-553.109.1.el8_10.x86_64}. This target belongs to a long-lived
enterprise kernel line in which vendor fixes and selected features are
backported while the reported upstream base version remains stable.
Consequently, the version number alone is not sufficient to determine whether
a hook, helper, program type, attachment mechanism, or verifier capability is
available \cite{redhat2025backporting}.

The design assumes that the monitored host exposes usable kernel
\gls{btf} information and grants the privileges required to load and attach the
selected programs. BTF and \gls{core} reduce dependence on compile-time kernel
structure layouts by allowing compatible field accesses to be relocated for
the deployed kernel. They do not provide facilities that the target lacks,
bypass permissions, or ensure that a program will be accepted by its verifier
\cite{linux_btf,linux_bpf_llvm_relocations,linux_libbpf_overview,
linux2026ebpfverifier}. Each required attachment and access path must therefore
be validated against the actual system.

This target-aware position deliberately limits the portability claim. The
architecture is CO-RE-assisted and is designed around the identified Rocky
Linux environment; it is not assumed to support every distribution or kernel
configuration unchanged. Where dynamic probes are used, symbol availability
and compatible function signatures remain additional assumptions. These
constraints motivate small kernel programs, explicit compatibility boundaries,
and a user-space layer that can report loading or attachment failures before
normal monitoring begins.

\subsection{Functional Requirements}
\label{subsec:functional-requirements}

Table~\ref{tab:functional-requirements} summarises the binding functional
requirements. They describe externally meaningful responsibilities rather than
individual commands or source-code components.

\begin{table}[H]
    \centering
    \small
    \renewcommand{\arraystretch}{1.18}
    \begin{tabularx}{\textwidth}{
        @{}
        >{\raggedright\arraybackslash}p{0.11\textwidth}
        >{\raggedright\arraybackslash}X
        @{}
    }
        \toprule
        \textbf{ID} & \textbf{Functional requirement} \\
        \midrule
        FR-01 &
        Collect a selectable set of process- and security-relevant host
        events. \\

        FR-02 &
        Resolve selected logical events into the producer and support probes
        required to create them, while excluding unrelated registered
        programs from loading and attachment. \\

        FR-03 &
        Transfer kernel records through the implemented perf-buffer channel
        and decode them into normalised typed events through an explicit
        shared contract. \\

        FR-04 &
        Load configurable policies that constrain the monitored event scope
        and may narrow attachment when their event allowlists are exact. \\

        FR-05 &
        Load detectors independently from policies and evaluate deterministic
        point conditions or bounded ordered collective sequences. \\

        FR-06 &
        Produce alerts that preserve the evidence responsible for a match and
        propagate available MITRE ATT\&CK metadata. \\

        FR-07 &
        Keep raw event presentation, alert presentation, and operational
        diagnostics as separate paths. \\

        FR-08 &
        Manage validation, startup, attachment, event consumption, observable
        failures, cancellation, and orderly resource cleanup. \\
        \bottomrule
    \end{tabularx}
    \caption{Functional requirements of the proposed system.}
    \label{tab:functional-requirements}
\end{table}

These requirements form a connected pipeline. FR-01 and FR-02 define what is
observed and which kernel programs are active. FR-03 establishes the
producer--consumer boundary. FR-04 and FR-05 separate monitoring scope from
detection logic, while FR-06 and FR-07 distinguish derived security results
from telemetry and diagnostics. FR-08 applies across the whole lifecycle.

The optional kernel-side UID guard is treated as one early reduction mechanism
within FR-02, not as a separate kernel policy language. Likewise, container
packaging is an implementation concern rather than a central functional
requirement of the architecture. A possible future orchestrated deployment
does not alter the host-local responsibilities defined here.

\subsection{Non-Functional and Operational Requirements}
\label{subsec:nonfunctional-operational-requirements}

The non-functional and operational requirements in
Table~\ref{tab:nonfunctional-requirements} constrain how the functional
pipeline is realised and evaluated.

\begin{table}[H]
    \centering
    \small
    \renewcommand{\arraystretch}{1.18}
    \begin{tabularx}{\textwidth}{
        @{}
        >{\raggedright\arraybackslash}p{0.11\textwidth}
        >{\raggedright\arraybackslash}X
        @{}
    }
        \toprule
        \textbf{ID} & \textbf{Non-functional or operational requirement} \\
        \midrule
        NFR-01 &
        Operate on the identified enterprise-kernel target with explicit
        compatibility assumptions and target validation. \\

        NFR-02 &
        Keep kernel execution and data extraction bounded and compatible with
        the target verifier. \\

        NFR-03 &
        Maintain a consistent kernel-to-user-space event contract and fail
        visibly on unknown or malformed records. \\

        NFR-04 &
        Keep detection deterministic, configuration-driven, and independent
        from event presentation. \\

        NFR-05 &
        Bound incomplete collective state through short windows and a
        per-detector key limit. \\

        NFR-06 &
        Expose transport loss, decoding failures, detector failures, and setup
        errors through appropriate diagnostic paths. \\

        NFR-07 &
        Evaluate a steady-state user-space CPU objective below five per cent of
        one core for explicitly defined representative profiles. \\
        \bottomrule
    \end{tabularx}
    \caption{Non-functional and operational requirements.}
    \label{tab:nonfunctional-requirements}
\end{table}

NFR-07 is deliberately phrased as an evaluation objective. It is not an
unconditional guarantee for every workload, event combination, or
instantaneous sample. Chapter~5 defines the profiles, warm-up periods, and
measurements used to assess it. Similarly, NFR-05 bounds retained correlation
keys and time, but the total memory demand also depends on the number and
complexity of configured detectors and on the size of retained evidence.

\section{Overall System Architecture}
\label{sec:overall-system-architecture}

The architecture is divided into a control path and a data path. The control
path validates configuration, prepares policies and detectors, derives the
effective event scope, and resolves logical events into physical programs
before loading. The data path begins when an attached kernel hook is reached
and ends with optional event or alert presentation. Operational diagnostics
observe both paths without becoming the representation used for events or
alerts.

\subsection{Kernel-Space and User-Space Responsibilities}
\label{subsec:kernel-userspace-responsibilities}

Kernel space is responsible for observing selected operations, extracting
bounded context and arguments, applying the optional exact-UID rejection,
constructing a logical record, and serialising it for submission. A limited
amount of temporary kernel state may join observations when an operation
requires information from both entry and return. This does not amount to a
kernel-resident policy or detector engine.

User space owns the more changeable and semantically rich responsibilities. A
Go runtime validates configuration, manages
\texttt{libbpf} through \texttt{libbpfgo}, resolves probe dependencies, receives
records, decodes typed arguments, normalises events, applies policy scope,
dispatches detectors, and produces presentation and diagnostic outputs
\cite{aquasecurity_libbpfgo}. Moving these responsibilities out of the kernel
limits verifier-sensitive logic and allows rule definitions to evolve without
recompiling each eBPF program.

The boundary is therefore selective rather than absolute. User space decides
which kernel programs are loaded and can configure a small kernel-side filter.
Kernel programs may retain short-lived construction state, but they do not own
the higher-level monitoring policy. Conversely, user-space analysis relies on
the correctness of the records produced in the kernel and cannot recover an
observation that was never submitted or was lost in transport.

\subsection{End-to-End Event Lifecycle}
\label{subsec:end-to-end-event-lifecycle}

Before collection begins, configuration and rule definitions are validated.
Where active policy rules provide an exact event allowlist, this scope may
narrow the requested event set. The probe registry then resolves public
logical events into the producer programs and internal dependencies needed to
construct them. Unrelated registered programs are excluded from the loading
and attachment path.

Once monitoring is active, a selected hook invokes its eBPF program. The
program may reject the event through the optional UID guard; otherwise it
captures bounded data, constructs the binary record, and submits it through the
perf buffer. The Go runtime decodes the record with the event registry, applies
the configured event and process filters, and asks the policy layer whether the
event belongs to the active scope. An admitted event may be presented as raw
telemetry and is then offered to relevant detectors. A successful detector
match produces an alert containing its source evidence.

Figure~\ref{fig:system-responsibility-flow} summarises these two paths. The
operational logging lane is separate because it reports lifecycle, attachment,
transport, decoding, and detector failures rather than carrying monitored
events.

\begin{figure}[H]
    \centering
    \small
    \fbox{%
        \begin{minipage}{0.94\textwidth}
            \textbf{Control path}\par
            Configuration and rule validation
            $\longrightarrow$ effective event scope
            $\longrightarrow$ probe and dependency resolution
            $\longrightarrow$ selective loading and attachment

            \medskip
            \hrule
            \medskip

            \begin{tabularx}{\textwidth}{
                >{\centering\arraybackslash}X
                >{\centering\arraybackslash}p{0.05\textwidth}
                >{\centering\arraybackslash}X
                >{\centering\arraybackslash}p{0.05\textwidth}
                >{\centering\arraybackslash}X
            }
                \textbf{Kernel space} &
                &
                \textbf{Go user space} &
                &
                \textbf{Presentation} \\
                \smallskip
                Selected hook &
                $\longrightarrow$ &
                Perf-buffer reader &
                $\longrightarrow$ &
                Raw event output \\
                eBPF capture &
                &
                Decoder and event registry &
                &
                Alert output \\
                Optional UID rejection &
                &
                Policy event scope &
                &
                Evidence and ATT\&CK metadata \\
                Binary serialisation &
                &
                Detector dispatch &
                &
                \\
                Perf-buffer submission &
                &
                &
                &
            \end{tabularx}

            \medskip
            \hrule
            \medskip

            \textbf{Operational diagnostics}\par
            A separate structured logging path reports setup, attachment,
            transport loss, decoding, detector, output, and cleanup
            diagnostics.
        \end{minipage}%
    }
    \caption{Responsibility boundary and end-to-end flow of the proposed
    monitoring system.}
    \label{fig:system-responsibility-flow}
\end{figure}

The diagram is intentionally not a guarantee of total event ordering. Perf
transport uses per-CPU buffering, and detector sequence order follows the order
in which records reach the user-space engine. Transport and ordering boundaries
are discussed further in Section~\ref{subsec:perf-buffer-failure-boundaries}.

\section{Kernel-Space Collection Design}
\label{sec:kernel-space-collection-design}

Kernel instrumentation must expose evidence that is both meaningful for the
logical event and obtainable on the target kernel. The design therefore does
not prescribe one universal attachment mechanism. It combines available
observation points while maintaining a stable event abstraction above them.

\subsection{Hook Selection and Logical Events}
\label{subsec:hook-selection-logical-events}

Four criteria guide hook selection. The first is the semantic position of the
observation: the hook must expose the request, resolved object, state
transition, or security mediation stage needed by the event. The second is
whether the result of the operation is required. An entry observation can
describe an attempt, but a return-side observation is needed when success or
failure changes the event's meaning.

The third criterion is availability on the concrete target. Tracepoints offer
explicit event interfaces when present, whereas dynamic probes extend coverage
at the cost of stronger dependence on symbols, signatures, and calling
conventions. The fourth criterion is expected cost and robustness.
High-frequency generic hooks can provide shared coverage but place filtering on
a hot path; a dedicated observation point is preferable when it provides the
required semantics reliably. These trade-offs follow from the properties of
the kernel mechanisms introduced in Chapter~2 rather than from an automatic
runtime choice.

The architecture distinguishes four related concepts. A physical hook is the
kernel execution point. An eBPF probe is the compiled program attached to that
point. A support probe is an internal program that preserves or supplies
context needed by another program. A logical event is the stable observation
exposed to the decoder, policies, detectors, and outputs.

These concepts are not necessarily one-to-one. A single logical operation may
require an entry probe to preserve request information and an exit probe to add
its result. Several operation variants may also converge on one logical event.
Conversely, selecting a public event may require internal support probes that
do not define independently selectable or decodable events. Keeping attachment
units separate from public event units prevents probe mechanics from becoming
part of the user-facing event contract.

\subsection{Early Event Selection and Filtering}
\label{subsec:early-event-selection-filtering}

Event selection begins with logical event names rather than individual eBPF
programs. The registry expands the selected events into producer and support
probes, after which unrelated registered programs are removed from the
autoload set. Only the resolved programs are presented to the normal loading,
verification, and attachment lifecycle. This design reduces active hooks and
avoids verifying registered programs that cannot contribute to the requested
event domain.

This mechanism is selective loading, not arbitrary elimination of every
unused section in the eBPF object. Its correctness depends on the registry
describing all public producers and their dependencies consistently. It also
does not remove the need for a later user-space gate: several physical
programs can contribute to one event, and policy or process constraints may
depend on information available only after decoding.

An optional exact-UID guard provides a second, narrower reduction. When
enabled, it rejects observations whose current effective UID differs from the
configured value before full context enrichment and serialisation. This saves
work for rejected observations, but intentionally supports only one equality
condition. It is not a range filter, a namespace or workload selector, or a
replacement for policy evaluation in user space.

Policies can influence early selection only when their active allow rules
produce an exact event scope. More specific policy constraints and suppression
decisions remain in user space. Detector definitions are separate consumers:
their declared event requirements are validated, but they are not
automatically converted into additional probe dependencies. A deployment must
therefore configure a policy or event scope that supplies the observations
required by its loaded detectors.

\section{Event Contract and Transport Design}
\label{sec:event-contract-transport-design}

The producer and consumer need a common representation that remains stable
across heterogeneous event sources. The design uses a compact binary contract
for transport and an explicit registry for semantic interpretation. These two
parts must evolve together: compact records reduce per-event overhead, but
they cannot be decoded safely when identifiers or schemas drift.

\subsection{Stable Binary Event Contract}
\label{subsec:stable-binary-event-contract}

Each public record contains a common event context followed by a sequence of
indexed, typed arguments. The common context provides the identity and
provenance needed across event families, including process, host, namespace,
and execution information. Event-specific arguments carry the values required
to interpret the observed operation. Their indices connect the encoded values
to the ordered schema held by the user-space event registry.

The record is compact rather than fully self-describing. The producer emits an
event identifier and indexed values, while the consumer supplies the event
name, expected argument types, and argument order. This separation reduces
repeated metadata in the transport path, but creates an explicit compatibility
obligation between the C producer, Go decoder, event registry, and probe
registry.

The contract must satisfy three properties. First, both sides must agree on the
common prefix. Second, every public event identifier must resolve to one
decoder schema. Third, argument indices and encodings must remain compatible
with that schema. An unknown identifier or structurally inconsistent record is
therefore treated as a contract failure. It is reported and dropped before
policy, detector, or presentation processing rather than being converted into
an untrusted partial event.

This design also separates structural and semantic processing. Binary decoding
establishes record validity and typed values. Registry lookup supplies the
logical event schema. Normalisation then presents one internal event model to
policy, detector, and output components. The concrete layout, field offsets,
buffer capacities, and compatibility corrections that enforce this contract
are implementation evidence and are deferred to Chapter~4.

\subsection{Perf-Buffer Transport and Failure Boundaries}
\label{subsec:perf-buffer-failure-boundaries}

The operational transport is a perf buffer backed by a perf-event array. This
mechanism is available on the target kernel and is supported by the selected
\texttt{libbpfgo} runtime stack. The BPF ring buffer discussed in Chapter~2 is
not implemented in the current event path. This is a compatibility-driven
engineering choice rather than a general claim that perf transport is
preferable for every monitoring system.

Perf transport provides asynchronous delivery from kernel producers to a
user-space consumer through finite per-CPU buffers
\cite{linux_perf_ring_buffer}. It can report that records were lost, but it
does not provide durable queuing, retransmission, or reconstruction of the
missing events. A submission failure in the eBPF program, a transport loss
notification, and a decoder failure after delivery are therefore distinct
conditions.

The system exposes these boundaries through diagnostics. Loss notifications
indicate that records were not delivered, while decoding errors indicate that
a delivered record could not be interpreted through the active contract.
Neither condition implies that the originating kernel operation failed to
occur. Output and detector failures arise later and likewise do not alter the
original operation.

Per-CPU delivery also limits ordering claims. A local sequence detector uses
the arrival order observed by the Go runtime; the architecture does not
reconstruct a total causal order across CPUs. Short correlation windows and
stable local identities reduce accidental associations, but ordering behaviour
under concurrent load remains an evaluation concern for Chapter~5.

\section{User-Space Runtime Design}
\label{sec:user-space-runtime-design}

The Go runtime coordinates kernel-facing resources and the higher-level event
pipeline. Its components are organised around validation before resource
acquisition, explicit ownership of the loaded runtime, and isolation of
event-local failures from the main processing loop.

\subsection{Loading, Attachment and Runtime Lifecycle}
\label{subsec:loading-attachment-runtime-lifecycle}

Startup proceeds through seven conceptual stages. Configuration and external
rule definitions are first normalised and validated. The eBPF object and BTF
source are then resolved. Next, policy information is used where possible to
derive the effective event scope, and logical events are expanded into the
required probes. The runtime creates the kernel object, excludes unrelated
registered programs from autoload, and asks the kernel to relocate, verify, and
load the remaining programs.

After a successful load, the optional filter configuration is written, the
perf-buffer consumer is started, and the selected programs are attached. The
event loop then waits for records, loss notifications, cancellation, or a
terminal runtime failure. This ordering prevents normal monitoring from
starting with a partially accepted configuration or an incomplete set of
attachments.

Setup failures are fatal because the requested observation domain cannot be
trusted when object loading, verification, transport initialisation, or
attachment fails. Event-local failures are handled differently: an invalid
record, detector error, or presentation error is reported and isolated so that
later records can still be processed.

Cancellation initiates the destruction of retained links and the closure of
transport, cgroup, and object resources in an ordered sequence. The
architecture promises that cleanup is attempted; it does not assume that every
deferred cleanup or logging error is propagated as the final process result.
This distinction is retained for the implementation discussion in Chapter~4.

\subsection{Decoding, Normalisation and Event Registration}
\label{subsec:decoding-normalisation-registration}

The decoder is the trust boundary for data received from the perf buffer. It
first establishes that the common context and declared argument records fit
within the delivered payload. It then resolves the event identifier through an
explicit registry and decodes each indexed argument according to the expected
type. Only a structurally valid and schema-compatible record becomes a
normalised event.

The event registry and the probe registry describe different aspects of the
same logical interface. The event registry defines identifiers, names, and
argument schemas used by decoding and configuration validation. The probe
registry defines which physical programs produce those events and which
support dependencies they require. Keeping these registries separate allows
each to represent its own concern, while consistency checks are required to
ensure that every public producer has a corresponding decoder contract.

Normalisation creates the representation consumed by later components in user
space. Policies and detectors receive typed event information rather than
parsing terminal text or JSON. Presentation can therefore change without
altering rule evaluation. Detectors can also preserve structured source
evidence instead of reconstructing it from rendered output.

An unknown event identifier or malformed argument record is logged and dropped.
The runtime does not fabricate an opaque event or pass a partially decoded
record to a policy or detector. This behaviour makes registry drift and
contract violations visible while allowing the monitoring loop to continue
with subsequent valid records.

\subsection{Event Output, Alert Output and Operational Logging}
\label{subsec:event-alert-logging}

The runtime maintains three distinct output responsibilities. Raw event output
presents admitted normalised telemetry. Alert output presents results derived
from detector matches. Structured operational logging reports the state and
failures of the monitoring infrastructure itself.

Raw events and alerts may be rendered independently. A presentation mode can
suppress raw telemetry while retaining collection, decoding, policy admission,
detector execution, and alert generation. This is an output decision and must
not be confused with early kernel filtering. In the combined mode, an admitted
raw event is presented before detector evaluation and any resulting alerts.

Operational diagnostics follow a separate path. They cover startup,
attachment, transport loss, decoding errors, detector failures, output
failures, and cleanup information. They neither serialise eBPF events nor
replace detector alerts. This separation allows machine-oriented telemetry and
security results to remain stable while logging verbosity or encoding changes
for operational needs.

Some failures can occur before the structured logger has been created, such as
invalid initial configuration or unreadable rule files. Those errors are
returned directly to the entry point. This does not change the three-way
architectural distinction; it identifies an outer startup boundary around the
normal runtime logging path.

\section{Policy and Detection Architecture}
\label{sec:policy-detection-architecture}

Policies and detectors are separate, configuration-driven stages over the
normalised event stream. This separation follows an established runtime
security pattern rather than constituting novelty by itself. Its role in the
prototype is to keep collection scope, security interpretation, and
presentation independently configurable within a bounded host-local design.

\subsection{Policy-Based Event Scope}
\label{subsec:policy-based-event-scope}

A policy determines whether a normalised event belongs to the active monitoring
scope. The implemented selectors use event identity and selected process
attributes. A matching suppression rule takes precedence over non-suppress
matches; when no policy is configured, the policy layer does not restrict the
event stream.

Some policy information can influence collection before monitoring starts.
When every active non-suppress rule supplies an explicit event allowlist, their
combined scope can narrow the event selection used by the probe registry.
Constraints that require decoded process information, together with
suppression decisions, remain in the user-space path. Policy evaluation is
therefore split between optional attach-time reduction of the event domain and
the authoritative admission decision over decoded events. It is not executed
as a complete kernel policy engine.

Policy definitions do not instantiate detectors, select detector identifiers,
or route detections by ATT\&CK metadata. Their mode and intent metadata also do
not create separate detector pipelines in the current runtime. Matched policy
names are not propagated automatically into emitted detector alerts. These
boundaries distinguish the implemented scope model from the broader policy
capabilities reviewed in Chapter~2.

\subsection{Point and Collective Detectors}
\label{subsec:point-collective-detectors}

Detector definitions are loaded independently from policies. Their event
references and conditions are validated before the runtime begins, and the
dispatcher indexes them by the events they consume. This avoids evaluating
every detector for every event and isolates a detector-local error from the
remaining detector set.

A point detector evaluates deterministic predicates over one normalised event.
When the predicates match, it creates an alert containing that source event.
A collective detector instead defines an ordered sequence of event conditions.
It retains a partial match until a later admitted event advances the sequence,
the sequence completes, or its local state is discarded. Neither model learns
a behavioural baseline or calculates a statistical anomaly score.

Detector event declarations are consumer requirements, not attachment
instructions. They are checked against the event registry, but they do not
automatically enlarge the collection scope. An operator must select compatible
events through the monitoring configuration or policies. A valid detector may
otherwise remain inactive because one of its required observations is not
produced.

Sequence order is the order in which events arrive at the user-space detector
engine. The design does not search backwards, reorder observations, or infer a
causal graph. This bounded model supports the short collective patterns defined
for the prototype while avoiding a claim of general complex-event processing
\cite{cugola2012processing,agrawal2008eventstreams}.

\subsection{Bounded Local Correlation}
\label{subsec:bounded-local-correlation}

Each collective detector defines how the events in its ordered sequence must
be related. The available local relationships distinguish a stable process, an
immediate parent-to-child transition, the same file resource, the same local
cgroup, or a conjunction of these dimensions. A partial match is retained only
within a short configurable interval and a bounded number of correlation keys.
When the sequence completes, the alert preserves the ordered events that
satisfied it; otherwise, the incomplete state expires, is replaced for the
same key, or is evicted.

Stable process correlation distinguishes separate process lifetimes even when
the operating system later reuses the same numeric process identifier.
Process-tree correlation remains deliberately shallow: a sequence may stay
within one process or move from a parent observation to its immediate child,
but it does not traverse arbitrary ancestry or construct a persistent process
graph.

Resource correlation identifies the same local file object rather than relying
on a pathname string. This permits related activity from different processes
to be joined while avoiding the assumption that an equal path always denotes
the same object. Cgroup correlation can similarly connect processes sharing a
local kernel cgroup identity. It does not attribute events to a Kubernetes pod,
workload, tenant, or cluster-wide identity.

A composite relationship requires every configured dimension to agree; it is
a conjunction rather than a fallback among alternative keys. User-session
correlation is rejected because the normalised event model does not expose a
stable session identity, and UID alone would merge unrelated activity.

Only one partial sequence is retained for a detector and correlation key. A
new first step for that key can replace the existing incomplete sequence.
State is bounded by time and by a per-detector key limit, although total memory
also depends on detector count, sequence length, and retained event payloads.
Expired state is reclaimed during relevant detector activity or an explicit
flush, not by a general periodic scheduler in the main runtime.

\subsection{Alert Evidence and MITRE ATT\&CK Metadata}
\label{subsec:alert-evidence-attack}

An alert is a derived detector result rather than an operational log message.
A point alert preserves the source event that satisfied its predicates, while
a collective alert preserves the ordered sequence accumulated by the detector.
This evidence makes the result traceable to the runtime observations on which
it is based and allows machine-oriented output to retain more detail than a
compact interactive summary.

Detector definitions may associate this evidence with MITRE ATT\&CK tactics,
techniques, data sources, and data components. The metadata is copied into the
resulting alert, providing a shared vocabulary through which the intended
security meaning can be communicated and organised
\cite{strom2020attack}. The mapping is part of detector configuration rather
than policy selection.

Validation establishes that configured tactic and technique identifiers follow
the expected syntactic form. It does not consult an authoritative ATT\&CK
catalogue at runtime, verify the semantic correctness of the mapping, or
demonstrate complete coverage. These limitations do not remove the value of
the metadata: preserving the mapping alongside the evidence supports
interpretation, traceability, and later coverage analysis. The final assessment
of representative detector mappings belongs to Chapter~5.

\section{Design Decisions and Scope Boundaries}
\label{sec:design-decisions-boundaries}

The architecture follows a series of deliberate trade-offs. Bounded
kernel-space collection is preferred to embedding complex, verifier-sensitive
analysis in each probe. The Go runtime owns policy and detector semantics so
that those definitions can evolve without recompiling the kernel programs.
Selective loading and attachment are preferred to activating every registered
probe, and an explicit shared event contract is preferred to ad hoc payload
interpretation.

Perf buffers are used because they are supported by the reference environment
and runtime stack; BPF ring-buffer transport remains a non-implemented
alternative. Policy scope is kept separate from detector logic, and detector
results are kept separate from operational logging. Collective detection uses
short, local, evidence-preserving sequences rather than persistent graphs or a
general-purpose event-processing system.

These decisions define the limits of the prototype. It does not implement
enforcement, prevention, machine-learning anomaly discovery, user-session
correlation, arbitrary process ancestry, cluster-wide correlation, or
Kubernetes workload attribution. It does not claim complete Linux event or
MITRE ATT\&CK coverage, universal kernel portability, or universal compliance
with the performance objective.

Container packaging supports later operational use but is an implementation
concern described in Chapter~4. A wider platform may eventually aggregate
alerts from several monitored hosts and add deployment context, but that
external capability is neither implemented nor evaluated as part of this
thesis. Within its stated boundary, the proposed system provides a coherent
path from selected kernel observations to normalised telemetry and
evidence-preserving local detections on the identified enterprise kernel.
